% =============================================================================
% Documentation technique — TiiBnTick Agency
% Version 2.0 — Analyse basée sur le code (frontend + BFF + docs backend)
% Compilation : pdflatex main.tex (×2)  |  latexmk -pdf main.tex
% =============================================================================
\documentclass[11pt,a4paper,openany]{report}

\usepackage[utf8]{inputenc}
\usepackage[T1]{fontenc}
\usepackage[french]{babel}
\usepackage{lmodern}
\usepackage{geometry}
\usepackage{graphicx}
\usepackage{hyperref}
\usepackage{booktabs}
\usepackage{longtable}
\usepackage{tabularx}
\usepackage{multirow}
\usepackage{enumitem}
\usepackage{listings}
\usepackage{xcolor}
\usepackage{fancyhdr}
\usepackage{titlesec}
\usepackage{microtype}
\usepackage{csquotes}
\usepackage{array}
\usepackage{caption}
\usepackage{needspace}

\geometry{margin=2.2cm, top=2.5cm, bottom=2.6cm}

\definecolor{brandorange}{HTML}{EA580C}
\definecolor{codebg}{HTML}{F5F5F4}
\definecolor{codeframe}{HTML}{D6D3D1}
\definecolor{ink}{HTML}{1C1917}

\hypersetup{
  colorlinks=true,
  linkcolor=brandorange,
  urlcolor=brandorange,
  citecolor=brandorange,
  pdftitle={TiiBnTick Agency — Documentation technique v2},
  pdfauthor={Équipe ingénierie TiiBnTick}
}

\lstset{
  backgroundcolor=\color{codebg},
  frame=single,
  rulecolor=\color{codeframe},
  basicstyle=\ttfamily\scriptsize,
  breaklines=true,
  columns=fullflexible,
  showstringspaces=false,
  tabsize=2,
  captionpos=b,
  aboveskip=0.8em,
  belowskip=0.8em
}

\pagestyle{fancy}
\fancyhf{}
\fancyhead[L]{\small\sffamily TiiBnTick Agency --- Doc.\ technique v2.0}
\fancyhead[R]{\small\thepage}
\fancyfoot[C]{\footnotesize\sffamily Usage interne --- ne pas diffuser hors équipe}
\renewcommand{\headrulewidth}{0.4pt}

\titleformat{\chapter}[display]
  {\normalfont\huge\bfseries\color{brandorange}}
  {\chaptertitlename\ \thechapter}{14pt}{\Huge}
\titleformat{\section}
  {\normalfont\Large\bfseries\color{ink}}{\thesection}{1em}{}
\titleformat{\subsection}
  {\normalfont\large\bfseries}{\thesubsection}{1em}{}

\newcommand{\pathref}[1]{\texttt{\detokenize{#1}}}
\newcommand{\api}[1]{\texttt{#1}}
\newcommand{\envvar}[1]{\texttt{\detokenize{#1}}}
\newcommand{\cls}[1]{\texttt{#1}}
\newcommand{\noteeng}[1]{\paragraph{Note d'ingénierie.} #1}

\begin{document}

\begin{titlepage}
  \centering
  \vspace*{1.8cm}
  {\sffamily\Large TiiBnTick}\\[0.35cm]
  {\sffamily\huge\bfseries Documentation technique}\\[0.35cm]
  {\sffamily\LARGE Agency --- Frontend, BFF \& intégrations}\\[1.2cm]
  \rule{\textwidth}{1.5pt}\\[0.5cm]
  {\large Analyse d'architecture et de conception à partir du code source}\\[0.25cm]
  {\large Plateforme de gestion d'agences de livraison}\\[0.5cm]
  \rule{\textwidth}{1.5pt}\\[1.4cm]
  \begin{tabular}{@{}rl@{}}
    \textbf{Version document} & 2.0 \\
    \textbf{Date} & \today \\
    \textbf{Périmètre code} & \pathref{frontend/} (runtime) + \pathref{backend/tnt-agency/} (rollback) \\
    \textbf{Upstream} & TiiBnTick Core (agency-registry) / Kernel (IAM via Core) \\
    \textbf{Audience} & Ingénieurs, architectes, DevOps, reviewers techniques \\
    \textbf{Statut} & Vivant --- à mettre à jour à chaque changement de contrat BFF \\
  \end{tabular}
  \vfill
  \begin{minipage}{0.92\textwidth}
    \small
    Ce document n'est \textbf{pas} un guide utilisateur (celui-ci est servi dans
    l'application sous \api{/guide}). Il décrit la stack, les frontières de
    responsabilité, les processus bout-en-bout, les contrats HTTP, les
    mécanismes d'auth, le mode offline livreur, et les dettes techniques
    constatées dans le dépôt.
  \end{minipage}
\end{titlepage}

\clearpage
\chapter*{Journal des versions}
\begin{tabular}{@{}llp{0.55\textwidth}@{}}
\toprule
\textbf{Ver.} & \textbf{Date} & \textbf{Contenu} \\
\midrule
1.0 & 2026-07-29 & Esquisse surfacique (rejetée --- trop générique) \\
2.0 & \today & Réécriture ingénierie : stack exacte, BFF, process, dettes \\
\bottomrule
\end{tabular}

\tableofcontents
\clearpage

\chapter{Objet, audience et conventions}

\section{Objet}

Documenter \textbf{ce qui est réellement implémenté} dans TiiBnTick Agency :
\begin{itemize}
  \item le frontend Next.js multi-portails ;
  \item le BFF embarqué \api{/api/agency/*} (chemin d'exécution principal) ;
  \item le BFF Java \pathref{backend/tnt-agency} (conservé en rollback) ;
  \item les contrats avec Core (agency-registry, auth, SSO, sync) et les satellites.
\end{itemize}

\section{Hors périmètre}

\begin{itemize}
  \item Guide utilisateur produit (\api{/guide} --- 17 sections, demos interactives).
  \item Spécification interne de Kernel ou de chaque module Core L0--L5.
  \item Landing écosystème \pathref{landing/} (hub multi-produits) --- mentionnée
        uniquement comme adjacent.
\end{itemize}

\section{Conventions de lecture}

\begin{description}[style=nextline]
  \item[Runtime actuel]
    Navigateur $\rightarrow$ Next BFF $\rightarrow$ Core. Affirmé dans
    \pathref{frontend/README.md}.
  \item[Rollback]
    BFF Java \texttt{tnt-agency} (:8081, base-path \api{/agency}). Même contrat
    conceptuel (agency-registry), déploiement distinct.
  \item[Écart documentaire]
    \pathref{ARCHITECTURE\_CORE\_KERNEL\_AGENCY.md} décrit encore une Agency avec
    DB locale et KernelBridge --- modèle historique. Le BFF Java actuel est
    \textbf{stateless} (R2DBC/Liquibase désactivés, \pathref{INTEGRATIONS.md}).
    Ce document priorise le \textbf{code runtime}.
\end{description}

\section{Glossaire opérationnel}

\begin{longtable}{@{}p{0.28\textwidth}p{0.64\textwidth}@{}}
\toprule
\textbf{Terme} & \textbf{Définition dans Agency} \\
\midrule
\endhead
Tenant & UUID plateforme Kernel/Core. Côté public :
  \envvar{NEXT\_PUBLIC\_TENANT\_ID}. Distinct de \texttt{agencyId}. \\
AgencyId & Identifiant d'une agence dans agency-registry. Provenance : QR, query
  string, session staff, réponse tracking. \\
Antenne / Branch & Site opérationnel rattaché à une agence. Portail
  \api{/branch/*}. \\
Hub relais & Point de dépôt/retrait. Occupancy + colis. \\
Mission & Unité de livraison. Statuts :
  \texttt{DRAFT|PENDING|ASSIGNED|IN\_TRANSIT|AT\_HUB|DELIVERED|FAILED|CANCELLED}
  (\pathref{lib/types.ts}). \\
Intake & Demande d'expédition / dépôt (QR public ou walk-in HQ). \\
POD & Proof of Delivery (photo, signature, OTP). \\
Platform Client & Couple \texttt{X-Client-Id}/\texttt{X-Api-Key} enregistré sur
  Core pour Agency. \\
Shared session & Cookie \texttt{tnt-shared-session} pour handoff SSO Yowyob. \\
\bottomrule
\end{longtable}

\section{Sources primaires (code)}

\begin{itemize}
  \item \pathref{frontend/package.json}, \pathref{middleware.ts},
        \pathref{lib/server/agency-bff.ts}, \pathref{lib/server/verify-auth.ts}
  \item \pathref{frontend/.env.example}, \pathref{lib/config.ts}, \pathref{lib/seo.ts}
  \item \pathref{backend/docs/AUTH\_BFF.md}, \pathref{ONBOARDING\_KERNEL\_FLOW.md},
        \pathref{ADR\_CORE\_FIRST\_MIGRATION.md}, \pathref{CORE\_CONSUMER\_CONTRACT.md}
  \item \pathref{backend/tnt-agency/INTEGRATIONS.md}
\end{itemize}

\chapter{Stack technologique}

\section{Frontend --- versions figées dans \texttt{package.json}}

Analyse de \pathref{frontend/package.json} (runtime Agency).

\begin{table}[h]
\centering
\begin{tabular}{@{}lll@{}}
\toprule
\textbf{Domaine} & \textbf{Technologie} & \textbf{Version / note} \\
\midrule
Framework & Next.js (App Router) & 15.5.20 \\
UI runtime & React / React DOM & 19.0.0 \\
Langage & TypeScript & \texttt{\^5} \\
CSS & Tailwind CSS + PostCSS + Autoprefixer & \texttt{\^3.4.17} \\
Formulaires & react-hook-form + @hookform/resolvers & \texttt{\^7.76} / \texttt{\^5.4} \\
Validation & Zod & \texttt{\^4.4.3} \\
JWT serveur & jose (JWKS) & \texttt{\^6.2.3} \\
Cartes & Leaflet + react-leaflet & \texttt{\^1.9.4} / \texttt{\^5.0} \\
Charts & recharts & \texttt{\^3.8.1} \\
Motion / icônes & framer-motion, lucide-react, clsx & --- \\
QR & qrcode.react & \texttt{\^4.2.0} \\
Realtime serveur & ws (pont STOMP/WS Core $\rightarrow$ SSE) & \texttt{\^8.21} \\
Lint & ESLint 9 + eslint-config-next & 15.5.20 \\
Tests & Vitest & \texttt{\^3.2.7} \\
Offline tests & fake-indexeddb & \texttt{\^6.0} \\
\bottomrule
\end{tabular}
\caption{Stack frontend Agency}
\end{table}

\subsection{Scripts npm}

\begin{lstlisting}
npm run dev | dev:turbo | build | start | lint | test
# test = "vitest run"
\end{lstlisting}

\subsection{Fichiers de tests présents}

\begin{itemize}
  \item \pathref{tests/agency-bff.test.ts}
  \item \pathref{tests/fleetman-bff.test.ts}
  \item \pathref{tests/offline-queue-idb.test.ts}
  \item \pathref{tests/offline-contracts.test.ts}
\end{itemize}

\noteeng{Couverture volontairement ciblée BFF/offline --- pas de suite E2E UI
dans ce dépôt au moment de l'audit.}

\subsection{Configuration Next}

\pathref{next.config.ts} : \texttt{reactStrictMode: true}. Route BFF :
\pathref{app/api/agency/[...path]/route.ts} avec \texttt{runtime: 'nodejs'} et
\texttt{dynamic: 'force-dynamic'}.

\section{BFF Java (rollback) --- profil technique}

\begin{itemize}
  \item Spring WebFlux, port \textbf{8081}, context path \api{/agency}
  \item Stateless : pas de schéma métier local (DataSource/Liquibase/R2DBC off)
  \item Clients WebClient vers Core agency-registry + auth/SSO/media/search/\ldots
  \item OpenAPI / controllers sous \texttt{…infrastructure.web}
\end{itemize}

\section{Upstream plateforme}

\begin{table}[h]
\centering
\begin{tabular}{@{}llp{0.42\textwidth}@{}}
\toprule
\textbf{Système} & \textbf{Host typique} & \textbf{Rôle pour Agency} \\
\midrule
TiiBnTick Core & tiibntick-core.yowyob.com & Source de vérité registry + proxy IAM \\
Kernel & kernel-core.yowyob.com & Émetteur JWT / comptes / orgs \\
Search & tiibntick-search.yowyob.com & Recherche \\
Trust & tiibntick-trust.yowyob.com & Preuves blockchain (pas auth) \\
FleetMan & fleetman.yowyob.com & Flotte avancée \\
\bottomrule
\end{tabular}
\caption{Systèmes amont}
\end{table}

\section{Choix d'architecture applicative (frontend)}

\begin{description}[style=nextline]
  \item[App Router]
    Séparation nette des portails par arborescence \pathref{app/}.
  \item[BFF same-origin]
    Le navigateur n'embarque jamais \envvar{TNT\_AGENCY\_API\_KEY}. Les appels
    métier partent vers \api{/api/agency}.
  \item[Session cookie HttpOnly]
    Après login, le JWT n'est plus exposé au JS (`accessToken` placeholder
    \texttt{http-only-session} côté client --- \pathref{lib/session.ts}).
  \item[Formulaires]
    RHF + Zod pour les créations métier (missions, policies, branches\ldots).
  \item[PWA livreur]
    Service worker + IndexedDB offline queue + sync pull/push/bootstrap.
\end{description}

\chapter{Architecture système}

\section{Règle d'or des couches}

\begin{lstlisting}[caption={Frontières Kernel / Core / Agency}]
Yowyob Kernel (IAM, org, RBAC, JWT RS256)
        ^
        |  UNIQUEMENT via Core (proxy /api/v1/auth/**, SSO, onboarding)
        |
TiiBnTick Core  (tnt-bootstrap)
  - agency-registry  /api/v1/tenants/{tid}/agency-registry/**
  - auth, media, KYC, realtime, SSO
        ^
        |  X-Client-Id + X-Api-Key (+ Bearer utilisateur)
        |
Agency BFF  (Next /api/agency  OU  Java /agency/v1)
        ^
        |  same-origin + cookie tnt-auth
        |
Navigateur multi-portails
\end{lstlisting}

Interdits documentés (\pathref{CORE\_CONSUMER\_CONTRACT.md}) :
\begin{itemize}
  \item Platform $\rightarrow$ Kernel en HTTP/JDBC/Kafka kernel direct ;
  \item secrets Core dans le bundle \texttt{NEXT\_PUBLIC\_*} ;
  \item confondre Trust (:8090) avec le service d'authentification.
\end{itemize}

\section{Dual BFF --- vérité opérationnelle}

\begin{table}[h]
\centering
\begin{tabularx}{\textwidth}{@{}lXX@{}}
\toprule
\textbf{Critère} & \textbf{Next \api{/api/agency}} & \textbf{Java \texttt{tnt-agency}} \\
\midrule
Statut runtime & \textbf{Principal} (README frontend) & Rollback / historique \\
Entrée HTTP & Route Handlers Node & WebFlux :8081 \api{/agency} \\
Auth & \pathref{agency-bff.ts} $\rightarrow$ Core & \cls{AuthBffController} $\rightarrow$ Core \\
Métier & Proxy registry + cas spéciaux & Controllers + \texttt{*ErpClient} \\
DB locale & Aucune & Aucune (stateless actuel) \\
\bottomrule
\end{tabularx}
\caption{Matrice dual BFF}
\end{table}

\noteeng{Toute évolution de contrat doit décider explicitement si le rollback
Java doit rester à parité. Aujourd'hui le username signup est volontairement
aligné (\texttt{deriveUsername}) pour éviter une divergence Kernel.}

\section{Schéma runtime Next}

\begin{lstlisting}
[Landing / Tarifs / Guide]   [HQ shell]   [/branch]   [/livreur PWA]   [/track]
            \________ LayoutController + middleware.ts ________/
                              |
                   fetch /api/agency/{path}
                              |
                 lib/server/agency-bff.ts
                              |
        +----------+----------+-----------+------------+
        |          |          |           |            |
     Core       Search    FleetMan   presence WS    SSO Core
   registry              (chiffr.)   -> SSE
\end{lstlisting}

\section{Identifiants multi-tenant}

\begin{enumerate}
  \item \textbf{Tenant plateforme} --- env public, pour portail anonyme.
  \item \textbf{AgencyId} --- résolu dynamiquement (QR \texttt{?agencyId=},
        session HQ après \api{auth/session}, tracking).
  \item Headers BFF$\rightarrow$Core : \texttt{x-tenant-id}, credentials
        plateforme, éventuellement Bearer utilisateur (cookie lu serveur).
\end{enumerate}

\section{Modèle organisationnel cible (docs)}

D'après \pathref{ARCHITECTURE\_CORE\_KERNEL\_AGENCY.md} : \\
\textit{1 entreprise (SARL) = 1 tenant = 1 organisation Kernel = 1 Agency TNT},
avec champs \texttt{kernel\_organization\_id}, \texttt{kernel\_business\_actor\_id}
après onboarding.

\section{ADR Core-first}

\pathref{ADR\_CORE\_FIRST\_MIGRATION.md} (Accepté, phases 1--6) : Core = source
de vérité logistique. Flags Java typiques :
\texttt{core-first-delivery|realtime|inventory|billing|incidents|media},
\texttt{strict-sync}. Frontend : \envvar{NEXT\_PUBLIC\_USE\_CORE\_REALTIME}.

Conséquence eng. : pas de « mode fantôme » métier si Core est down --- fail
visible plutôt que cohérence locale mensongère.

\chapter{BFF Next.js --- conception détaillée}

\section{Point d'entrée}

\begin{itemize}
  \item Route : \pathref{app/api/agency/[...path]/route.ts}
  \item Handler : \pathref{lib/server/agency-bff.ts} --- \texttt{handleAgencyBff}
  \item Auth JWT : \pathref{lib/server/verify-auth.ts}
  \item FleetMan : \pathref{lib/server/fleetman-bff.ts}
  \item Présence : \pathref{lib/server/presence-stream.ts}
\end{itemize}

\section{Pipeline de traitement d'une requête}

\begin{enumerate}
  \item \texttt{normalizeBffPath} --- refuse \texttt{..}, schemes, \texttt{\textbackslash 0}
        (anti path-traversal).
  \item Cas \api{auth/logout} --- clear cookies.
  \item Si route non publique $\rightarrow$ \texttt{verifyRequestToken}
        (Bearer ou cookie \texttt{tnt-auth}).
  \item Branches dédiées : auth, sessions enrichies, media/KYC, Yowyob launch,
        presence, search, fleetman, sync, SSE tracking/missions.
  \item \textbf{Fallback} : proxy Agency Registry \\
        \api{\{CORE\}/api/v1/tenants/\{tenantId\}/agency-registry/\{path\}}
\end{enumerate}

\section{Routes publiques BFF (\texttt{isPublicRequest})}

Sans JWT :
\begin{itemize}
  \item \texttt{POST} \api{auth/login}, \api{auth/signup}, \api{auth/login/mfa/confirm}
  \item \texttt{GET} \api{tracking/\{code\}} et \api{tracking/\{code\}/stream}
  \item \texttt{GET} \api{intake-requests/\{id\}}
  \item \texttt{GET} \api{agencies/\{id\}/relay-hubs|intake-context}
  \item \texttt{POST} \api{intake-requests}, \api{agencies/\{id\}/drop-off|claims}
  \item \texttt{POST} \api{hub-handoffs/\{id\}/confirm-client}
\end{itemize}

Sessions enrichies (JWT) : \api{auth/session}, \api{auth/branch/session},
\api{auth/livreur/session}, \api{auth/hub/session}.

Rewrites / handlers dédiés :
\begin{itemize}
  \item \texttt{hubs/*/parcels} POST $\rightarrow$ \texttt{\ldots/deposit}
  \item \texttt{hubs/expired} $\rightarrow$ \texttt{\ldots/process} + mapping \texttt{\{processed\}}
  \item withdraw par recordId $\rightarrow$ \texttt{hub-parcels/\{tracking\}/withdraw}
  \item \texttt{PATCH deliverers/\{id\}/location} $\rightarrow$ persist + \texttt{realtime/gps/ping}
  \item \texttt{GET admin/agencies} $\rightarrow$ page \texttt{\{items,page,size,total\}}
\end{itemize}

Toute nouvelle API anonyme doit être ajoutée ici \textbf{et} dans le middleware
edge si elle a une page associée.

\section{Envelope de réponse}

Helpers \texttt{success()} / \texttt{error()} :

\begin{lstlisting}
{
  "status": "SUCCESS" | "ERROR",
  "data": { ... } | null,
  "error": { "code": "...", "message": "..." } | null,
  "timestamp": "ISO-8601"
}
\end{lstlisting}

Client : \pathref{lib/api/envelope.ts} (\texttt{unwrapApiData}),
\pathref{lib/api/client.ts}, messages FR \pathref{lib/errors.ts}.

\noteeng{Hétérogénéité connue : certaines réponses session/media sont en JSON
plat (hors envelope). \texttt{unwrapApiData} tolère les deux formes --- dette
à résorber pour homogénéiser les contrats.}

\section{Mapping auth BFF $\rightarrow$ Core}

\begin{table}[h]
\centering
\begin{tabular}{@{}ll@{}}
\toprule
\textbf{BFF} & \textbf{Core} \\
\midrule
\api{POST \ldots/auth/login} & \api{POST /api/v1/auth/login} \\
\api{POST \ldots/auth/signup} & \api{POST /api/v1/auth/sign-up} \\
\api{POST \ldots/auth/login/mfa/confirm} & \api{POST /api/v1/auth/login/mfa/confirm} \\
\bottomrule
\end{tabular}
\end{table}

Signup force côté BFF : \texttt{accountType: BUSINESS},
\texttt{businessType: LOGISTICS}, username via \texttt{deriveUsername(email)}
(aligné BFF Java).

Outcomes documentés : \texttt{AUTHENTICATED} | \texttt{MFA\_REQUIRED} |
\texttt{EMAIL\_VERIFICATION\_REQUIRED}.

\section{Credentials serveur}

Injectés uniquement serveur :
\begin{itemize}
  \item \envvar{TNT\_CORE\_BASE\_URL}, \envvar{TNT\_CORE\_WS\_URL}
  \item \envvar{TNT\_AGENCY\_CLIENT\_ID}, \envvar{TNT\_AGENCY\_API\_KEY}
  \item Alias legacy acceptés : \texttt{CORE\_BASE\_URL}, \texttt{CORE\_CLIENT\_ID},
        \texttt{YOWAUTH0\_*}
\end{itemize}

\section{BFF Java --- cartographie controllers (rollback)}

Préfixe effectif \api{/agency/v1/…} (échantillon) :
\cls{AuthBffController}, \cls{OnboardingController},
\cls{OnboardingAdminController}, \cls{AgencyController},
\cls{MissionController}, \cls{HubController}, \cls{FleetController},
\cls{FleetManController}, \cls{BillingController}, \cls{IntakeController},
\cls{ClientController}, \cls{ComplianceController}, \cls{SyncController},
\cls{YowyobIntegrationController}, streams SSE.

Clients ERP : package \texttt{agencyerp/*ErpClient} sur
\api{/api/v1/tenants/\{tid\}/agency-registry}.

\chapter{Portails, middleware et authentification}

\section{Matrice des portails}

\begin{longtable}{@{}p{0.18\textwidth}p{0.28\textwidth}p{0.22\textwidth}p{0.22\textwidth}@{}}
\toprule
\textbf{Portail} & \textbf{Routes} & \textbf{Shell} & \textbf{Auth} \\
\midrule
\endhead
Marketing & \api{/}, \api{/tarifs}, \api{/pricing}$\rightarrow$tarifs & LandingChrome & Public edge \\
Guide & \api{/guide}, \api{/guide/[slug]} & GuideShell & Public \\
HQ & \api{/dashboard}, missions, flotte, hubs, staff, billing, accueil, litiges, messages, settings, profile & Sidebar+Header & Cookie + AuthGuard \\
Register & \api{/register}, \api{/pending} & Standalone & Public \\
Antenne & \api{/branch/*} & BranchAppShell & Public edge + auth portail \\
Livreur & \api{/livreur/*} & LivreurNav PWA & Public edge + auth portail \\
Track & \api{/track}, deposit, messages & Standalone & Public \\
Admin & \api{/admin/*} & AdminShell & JWT + rôle plateforme \\
\bottomrule
\end{longtable}

\section{LayoutController}

\pathref{components/LayoutController.tsx} :
\begin{itemize}
  \item Préfixes standalone : login, register, pending, track, livreur, branch,
        admin, guide, tarifs, pricing, et \api{/}.
  \item Sinon : \cls{ThemeProvider} + \cls{AuthGuard} + Sidebar + Header +
        \cls{PageTransition}.
\end{itemize}

\section{Middleware edge (\pathref{middleware.ts})}

\begin{enumerate}
  \item Ignore \api{\_next}, \api{/api}, fichiers statiques.
  \item Zone admin : hors \api{/admin/login}, exige JWT + rôle
        (\texttt{TNT\_ADMIN} / platform admin via claims).
  \item Routes publiques listées (dont \api{/tarifs}).
  \item Sinon : cookie \texttt{tnt-auth} vérifié via JWKS ; sinon redirect
        \api{/login?from=…}.
  \item Headers : \texttt{X-Robots-Tag} hors SEO public, nosniff, referrer-policy.
\end{enumerate}

\noteeng{Ajouter une page marketing sans mettre à jour le middleware provoque
une fausse « page de connexion » --- déjà observé pour le pricing avant
whitelist.}

\section{Cookies et session navigateur}

\begin{table}[h]
\centering
\begin{tabularx}{\textwidth}{@{}lXX@{}}
\toprule
\textbf{Artefact} & \textbf{Rôle} & \textbf{Détail} \\
\midrule
\texttt{tnt-auth} & JWT HttpOnly & sameSite=lax, secure en prod, maxAge 24h \\
\texttt{tnt-shared-session} & Handoff SSO Yowyob & Posé au login si fourni Core \\
localStorage & Métadonnées UX & tenant, agencyId, user, agencyStatus/active \\
\bottomrule
\end{tabularx}
\end{table}

Le client \textbf{ne lit plus} le JWT (\pathref{lib/session.ts} :
\texttt{getAuthToken() → null}). Les \texttt{fetch} same-origin envoient le
cookie automatiquement.

\section{Vérification JWT}

\pathref{lib/server/verify-auth.ts} : bibliothèque \textbf{jose}, JWKS distant
(\envvar{JWT\_ISSUER\_URI}, \envvar{JWT\_JWK\_SET\_URI}). Source token :
header \texttt{Authorization: Bearer} \textbf{ou} cookie.

\section{Enrichissement session HQ}

Après login authentifié :
\begin{enumerate}
  \item \api{GET /api/agency/auth/session}
  \item Résolution agency via registry (application \texttt{me}, agence, fallback email)
  \item Persistance : \texttt{tnt-agency-id}, \texttt{tnt-agency-status},
        \texttt{tnt-agency-active}, user meta (\pathref{authService.ts})
\end{enumerate}

Sessions antenne / livreur : endpoints dédiés
\api{auth/branch/session}, \api{auth/livreur/session} (résolution staff/deliverer
côté BFF --- coût N+1 possible, dette connue).

\section{AuthGuard client (HQ uniquement)}

\pathref{components/AuthGuard.tsx} :
\begin{itemize}
  \item Non authentifié $\rightarrow$ \api{/login} (ou admin login)
  \item Statut \texttt{REJECTED} $\rightarrow$ pending rejected
  \item Non \texttt{ACTIVE} $\rightarrow$ \api{/pending}
\end{itemize}

\section{MFA}

\begin{lstlisting}
POST login -> 202 / CONFIRM_MFA + mfaToken
UI MfaChallengeForm -> POST auth/login/mfa/confirm { mfaToken, code }
\end{lstlisting}

Services : \pathref{yowauthService.ts}, \pathref{useLoginWithMfa},
\pathref{components/auth/MfaChallengeForm.tsx}.

\chapter{Processus métier bout-en-bout}

\section{P1 --- Inscription entreprise et activation}

\subsection{Étapes produit}

\begin{enumerate}
  \item Création compte IAM (signup BFF $\rightarrow$ Core $\rightarrow$ Kernel).
  \item Formulaire multi-étapes \pathref{app/register/page.tsx}
        (identité, KYC docs, settings, credentials).
  \item \texttt{registerService.submitApplication} $\rightarrow$
        \api{/onboarding/applications}.
  \item Complétion identité Kernel :
        \texttt{completeKernelIdentity} / \texttt{ensureKernelIdentity}.
  \item Attente \api{/pending} (\texttt{getMyApplication}).
  \item Admin TNT : \api{/admin/onboarding/[id]} --- approve/reject
        (\pathref{onboardingAdminService.ts}).
\end{enumerate}

\subsection{Contrainte Kernel critique}

\pathref{ONBOARDING\_KERNEL\_FLOW.md} : l'endpoint actors/onboarding est lié au
\texttt{sub} du JWT. La phase 1 \textbf{doit} s'exécuter avec le JWT
\textbf{candidat}. Approuver avec le JWT admin sans phase 1 associée crée le
mauvais BusinessActor.

\subsection{Phase admin (orchestration Core)}

Approve $\rightarrow$ organisation Kernel + rôles TNT + options plateforme
(trust ON) + rôle manager. Prérequis : \texttt{kernelIdentityReady=true}.

SSO post-activation (HRM / KSM\ldots) exige \texttt{kernelOrganizationId}.

\begin{lstlisting}[caption={Séquence onboarding (simplifiée)}]
Candidat -> BFF signup -> Core auth/sign-up -> Kernel
Candidat -> BFF onboarding/applications (+ Bearer candidat)
Candidat -> BFF .../kernel-identity
Admin    -> BFF admin/.../approve -> Core onboarding approve
         -> Agency ACTIVE + org Kernel liée
\end{lstlisting}

\section{P2 --- Intake client et création de mission}

\subsection{Canal QR public}

\begin{enumerate}
  \item HQ génère URL : \texttt{intakeDepositUrl(agencyId, branchId)}
        $\rightarrow$ \api{/track/deposit?agencyId=\&branchId=}.
  \item Client charge contexte : \api{GET agencies/\{id\}/intake-context}.
  \item Submit : \api{POST intake-requests}.
  \item Agent HQ : file \api{/accueil/demandes} --- approve/reject.
  \item Walk-in direct : \api{/accueil} via \texttt{intakeService.walkIn}.
\end{enumerate}

\subsection{Canal suivi}

\api{/track} : recherche code, stream SSE optionnel, claims / drop-off / hubs
relais (\pathref{trackingService.ts}, \pathref{complianceService.ts}).

\section{P3 --- Lifecycle mission}

Statuts (\pathref{lib/types.ts}) :
\texttt{DRAFT → PENDING → ASSIGNED → IN\_TRANSIT → AT\_HUB → DELIVERED}
(ou \texttt{FAILED}/\texttt{CANCELLED}).

\begin{description}[style=nextline]
  \item[HQ / Antenne]
    create, assign/reassign, cancel/reschedule, deposit-hub, withdraw
    (\pathref{missionService.ts}, services branch).
  \item[Livreur]
    pickup, deliver (+ POD), anomaly ; si offline $\rightarrow$ queue IndexedDB
    puis sync (\pathref{livreurMissionService}, \pathref{offlineQueue.ts}).
\end{description}

\section{P4 --- Facturation et commissions}

\begin{itemize}
  \item Politiques tarifaires actives (devise XAF typique).
  \item Estimation à la création de mission.
  \item Factures / revenus après clôture.
  \item Commissions livreur (portail gains + HQ personnel).
\end{itemize}
Services : \pathref{billingService.ts}, \pathref{livreurCommissionService.ts}.

\section{P5 --- Flotte et FleetMan}

\begin{enumerate}
  \item CRUD véhicules Agency (registry).
  \item Pont FleetMan (managers HQ) :
        \api{agencies/\{id\}/fleetman/\{status|connect|launch|sync\}}.
  \item Tokens stockés chiffrés AES
        (\envvar{TNT\_FLEETMAN\_TOKEN\_ENCRYPTION\_KEY}) ;
        plaintext interdit en prod.
\end{enumerate}

\section{P6 --- Hubs et Trust}

Transitions hub / POD peuvent ancrer des preuves côté Trust \textbf{via Core}
après activation. Le frontend expose
\envvar{NEXT\_PUBLIC\_TNT\_TRUST\_BASE\_URL} mais n'implémente pas (audit) un
client Trust autonome --- la vérité passe par les APIs agency/Core.

\chapter{Domaines métier et services client}

\section{Cartographie \pathref{lib/services/}}

\begin{longtable}{@{}p{0.34\textwidth}p{0.58\textwidth}@{}}
\toprule
\textbf{Service} & \textbf{Responsabilité} \\
\midrule
\endhead
\pathref{authService.ts} & Session HQ/admin, persist localStorage \\
\pathref{yowauthService.ts} / couche auth & login/signup/MFA vers BFF \\
\pathref{registerService.ts} & Candidature onboarding \\
\pathref{onboardingAdminService.ts} & Approve/reject admin \\
\pathref{agencyService.ts} & Profil / settings agence \\
\pathref{branchService.ts} & Antennes HQ \\
\pathref{staffService.ts} & Managers, livreurs, contrats, freelancers \\
\pathref{fleetService.ts} / \pathref{fleetManService.ts} & Véhicules + pont FleetMan \\
\pathref{missionService.ts} & Missions HQ \\
\pathref{hubService.ts} & Hubs + colis \\
\pathref{billingService.ts} & Policies, invoices, estimate \\
\pathref{agencyOperationsService.ts} & Agrégats ops HQ \\
\pathref{analyticsService.ts} & KPI / reports \\
\pathref{notificationService.ts} & Inbox notifications \\
\pathref{searchService.ts} & Recherche (proxy Search) \\
\pathref{intakeService.ts} & Intake / walk-in / approve \\
\pathref{trackingService.ts} & Suivi public \\
\pathref{complianceService.ts} & Litiges / claims \\
\pathref{mediaService.ts} & Upload media / KYC \\
\pathref{adminService.ts} & Admin agences \\
\pathref{yowyobIntegrationService.ts} & Launch SSO apps \\
\pathref{branchAuthService.ts} + portal/ops & Portail antenne \\
\pathref{livreurAuthService.ts} + mission/hub/location/media/commission & Portail livreur \\
\bottomrule
\end{longtable}

\section{Couche API client}

\begin{itemize}
  \item \pathref{lib/api/client.ts} --- fetch same-origin, gestion envelope/HTTP
  \item \pathref{lib/api/envelope.ts} --- unwrap
  \item \pathref{lib/api/mappers.ts} + \pathref{dto.ts} --- DTO$\rightarrow$modèle UI
  \item \pathref{lib/errors.ts} + toast --- UX erreur non technique
\end{itemize}

\noteeng{\pathref{dto.ts} porte encore le commentaire « réponses brutes du
backend tnt-agency » --- nomenclature héritée alors que l'upstream runtime est
Core agency-registry.}

\section{Patterns UI métier}

\begin{description}[style=nextline]
  \item[Drawers]
    \pathref{components/forms/Drawer.tsx} + *DetailDrawer (mission, contrat,
    commission, freelancer, invoice, branch\ldots).
  \item[Forms]
    Create*Form sous \pathref{components/forms/} (RHF+Zod).
  \item[Action queues]
    \pathref{lib/agency/actionQueue.ts}, \pathref{lib/branch/actionQueue.ts}
    pour files d'actions ops.
  \item[Hooks data]
    \texttt{useAgencyOperationalData}, \texttt{useBranchOperationalData},
    positions GPS live.
  \item[Toast]
    \pathref{contexts/ToastContext.tsx} (4s auto-dismiss).
\end{description}

\section{Surfaces encore immatures (constat code)}

\begin{itemize}
  \item Messagerie : pages présentes + indicateurs « coming soon » /
        \cls{ComingSoonBadge}.
  \item Trust : URL configurée, pas de SDK client dédié.
\end{itemize}

\chapter{Intégrations externes}

\section{Tableau de synthèse}

\begin{longtable}{@{}p{0.18\textwidth}p{0.32\textwidth}p{0.42\textwidth}@{}}
\toprule
\textbf{Système} & \textbf{Mécanisme} & \textbf{Fichiers / env} \\
\midrule
\endhead
Core Registry & Proxy HTTP BFF & \pathref{agency-bff.ts}, \envvar{TNT\_CORE\_*} \\
Core Auth & Proxy /api/v1/auth & Platform Client ID/Key \\
Core Media/KYC & Routes BFF dédiées & \pathref{mediaService.ts} \\
Core Realtime & WS serveur $\rightarrow$ SSE client & \pathref{presence-stream.ts},
  \pathref{coreRealtime.ts} \\
Core SSO & /api/v1/sso/yowyob/launch & shared-session cookie \\
Search & Proxy \envvar{TNT\_SEARCH\_BASE\_URL} & échec $\rightarrow$ hits vides \\
FleetMan & status/connect/launch/sync & \pathref{fleetman-bff.ts}, AES tokens \\
Trust & Preuves via Core ; URL publique & \envvar{NEXT\_PUBLIC\_TNT\_TRUST\_BASE\_URL} \\
Yowyob apps & HRM, ACCOUNTING, BILLING, CASHIER\ldots &
  \pathref{yowyobIntegrationService.ts} \\
Maps/WhatsApp & Liens UI only & \texttt{NEXT\_PUBLIC\_MAPS\_*}, \texttt{WHATSAPP\_*} \\
Site global & CGU/privacy/cookies/social & \pathref{lib/config.ts} \texttt{GLOBAL\_LINKS} \\
\bottomrule
\end{longtable}

\section{Platform Client Core}

Avant tout appel prod, Agency doit disposer d'un client plateforme créé via
\api{/api/v1/admin/platform-clients} (rôle \texttt{TNT\_ADMIN}) puis d'une
API key (\texttt{plaintextSecret} une seule fois). Scopes : \texttt{*} ou
\texttt{AUTH:*}, \texttt{SSO:*}, \texttt{ONBOARDING:*}, etc. (guide
\pathref{core/docs/auth/platform-client-onboarding-guide.md}).

\section{SSO Yowyob --- process}

\begin{lstlisting}
UI YowyobLaunchCard
  -> GET /api/agency/integrations/yowyob/launch?app=HRM|...
  -> BFF ajoute shared session + tenant + agency (+ branch)
  -> Core /api/v1/sso/yowyob/launch
  -> redirect app Yowyob /sso/callback?token&organizationId&scope
\end{lstlisting}

Prérequis : onboarding approuvé (\texttt{kernelOrganizationId}). Rôles autorisés
selon l'app (managers, accountant, branch manager scopé).

\section{FleetMan --- particularités}

\begin{itemize}
  \item Pas de SSO Yowyob pour FleetMan (auth propre / password user --- commentaire code).
  \item Restriction fréquente aux managers HQ dans le BFF.
  \item Chiffrement tokens obligatoire hors local.
\end{itemize}

\section{Règles d'évolution}

\begin{enumerate}
  \item Nouvel externe = client \textbf{uniquement} dans \pathref{lib/server/*}.
  \item Documenter env dans \pathref{.env.example}.
  \item Mettre à jour ce chapitre + tests Vitest si contrat BFF change.
\end{enumerate}

\chapter{Temps réel et mode hors-ligne livreur}

\section{Temps réel --- deux chemins}

\subsection{Chemin par défaut (Core)}

\begin{itemize}
  \item Flag : \envvar{NEXT\_PUBLIC\_USE\_CORE\_REALTIME=true}
  \item Client : \texttt{EventSource} vers
        \api{/api/agency/realtime/presence}
  \item Serveur : \pathref{presence-stream.ts} ouvre WS/STOMP vers Core
        (\envvar{TNT\_CORE\_WS\_URL}) et relaie en SSE
  \item Impl client : \pathref{lib/coreRealtime.ts}, variantes branch
        \pathref{branchRealtime.ts}
\end{itemize}

\subsection{Chemin legacy}

Si flag \texttt{false} : WebSocket legacy branché sur la base API
(\pathref{lib/realtime.ts}). Coexistence = dette de simplification.

\subsection{SSE métier}

Streams tracking et événements mission exposés par le BFF (routes dédiées dans
\pathref{agency-bff.ts}).

\section{Offline livreur --- conception}

\subsection{Objectifs}

Continuer pickup / deliver / anomaly / deposit-hub / location sans réseau, puis
réconcilier.

\subsection{Composants}

\begin{table}[h]
\centering
\begin{tabularx}{\textwidth}{@{}lX@{}}
\toprule
\textbf{Module} & \textbf{Rôle} \\
\midrule
\pathref{lib/livreur/offlineQueue.ts} & File IndexedDB \texttt{tnt-livreur-offline} \\
\pathref{lib/livreur/syncService.ts} & pull / push / bootstrap $\rightarrow$
  \api{agencies/\{id\}/sync/*} \\
\pathref{lib/livreur/syncConflicts.ts} & Résolution conflits \\
\pathref{lib/offline/readModel.ts} & Cache read-model local \\
UI & \cls{OfflineSyncBanner}, SW (\cls{RegisterServiceWorker}), install PWA \\
Tests & \pathref{tests/offline-queue-idb.test.ts},
  \pathref{tests/offline-contracts.test.ts} \\
\bottomrule
\end{tabularx}
\end{table}

\subsection{Types d'opérations queue}

Pickup, deliver (POD), anomaly, deposit-hub, location --- alignés sur les
transitions livreur online.

\subsection{Exigences non fonctionnelles}

\begin{enumerate}
  \item Idempotence autant que le contrat Core le permet.
  \item Horodatage client + serveur pour départager.
  \item Aucune sync hors session livreur authentifiée (cookie).
  \item Toute évolution de schéma IndexedDB = migration versionnée + test.
\end{enumerate}

\begin{lstlisting}[caption={Boucle sync offline}]
Action terrain (offline)
  -> enqueue IndexedDB
  -> retour réseau
  -> syncService push
  -> Core sync via BFF
  -> pull/bootstrap refresh read-model
  -> UI retire / marque conflits
\end{lstlisting}

\chapter{Déploiement, configuration et SEO}

\section{Prérequis runtime}

\begin{itemize}
  \item Node.js $\geq$ 20, npm $\geq$ 9
  \item Platform Client Core valide
  \item JWKS Kernel joignable depuis le process Next
  \item Accès réseau Core / Search / FleetMan selon features activées
\end{itemize}

\section{Commandes}

\begin{lstlisting}
cd frontend
cp .env.example .env.local   # renseigner secrets
npm install
npm run dev                  # local
npm run build && npm run start
npm test
\end{lstlisting}

\section{Variables --- contrat \pathref{.env.example}}

\subsection{Publiques (build-time)}

\begin{itemize}
  \item \envvar{NEXT\_PUBLIC\_API\_BASE\_URL} /
        \envvar{NEXT\_PUBLIC\_AGENCY\_PUBLIC\_BASE\_URL} $\rightarrow$ \api{/api/agency}
  \item \envvar{NEXT\_PUBLIC\_AGENCY\_FRONTEND\_URL} --- SEO, QR absolus
  \item \envvar{NEXT\_PUBLIC\_TENANT\_ID} --- tenant plateforme public
  \item Carte : \texttt{MAP\_TILE\_URL}, centre Douala par défaut, attribution OSM
  \item \envvar{NEXT\_PUBLIC\_USE\_CORE\_REALTIME}
  \item Trust URL, site global, légal, réseaux sociaux
\end{itemize}

\subsection{Serveur (runtime secrets)}

\begin{itemize}
  \item Core : base URL, WS, client id, api key
  \item JWT : issuer + JWKS
  \item Search base URL
  \item FleetMan : API/UI, enabled, encryption key, allow plaintext=false prod
\end{itemize}

\noteeng{Changer une variable \texttt{NEXT\_PUBLIC\_*} impose un rebuild.
Les secrets ne doivent jamais utiliser ce préfixe.}

\section{SEO}

\begin{itemize}
  \item \pathref{lib/seo.ts} --- \texttt{buildPageMetadata},
        \texttt{PUBLIC\_SEO\_ROUTES}, \texttt{isSeoPublicPath}
  \item \pathref{app/sitemap.ts} --- routes publiques + toutes sections guide
  \item \pathref{app/robots.ts} --- allow marketing/track/login ; disallow
        \api{/api/}, HQ ops, \api{/livreur/}, \api{/branch/}
  \item JSON-LD landing : \pathref{components/seo/JsonLd.tsx}
\end{itemize}

\section{Checklist mise en production}

\begin{enumerate}[leftmargin=*]
  \item Secrets Core/FleetMan renseignés ; plaintext FleetMan off.
  \item \envvar{NEXT\_PUBLIC\_TENANT\_ID} = UUID tenant plateforme réel.
  \item Middleware : toutes pages publiques whitelistées.
  \item Smoke : login HQ MFA, register, approve admin, track QR, mission
        assignée, POD livreur, sync offline simulée.
  \item Vérifier SSO launch si HRM/KSM activés (org Kernel liée).
  \item \texttt{npm run build} + \texttt{npm test} verts.
\end{enumerate}

\section{Hosts de référence}

\begin{tabular}{@{}ll@{}}
\toprule
Agency UI & \texttt{https://tiibntick-agency.yowyob.com} \\
Core & \texttt{https://tiibntick-core.yowyob.com} \\
Kernel & \texttt{https://kernel-core.yowyob.com} \\
Entry / légal & \texttt{https://tiibntick.yowyob.com} \\
\bottomrule
\end{tabular}

\chapter{Sécurité, opérations et dette technique}

\section{Contrôles de sécurité implémentés}

\begin{itemize}
  \item JWT validé JWKS (edge + BFF).
  \item Cookie session HttpOnly ; pas de Bearer dans localStorage.
  \item Path normalization anti-traversal sur le BFF.
  \item Séparation Platform Client / RBAC utilisateur.
  \item Chiffrement tokens FleetMan.
  \item Headers basiques (nosniff, referrer, robots).
  \item Rôles admin plateforme sur \api{/admin/*}.
\end{itemize}

\section{Runbook --- symptômes fréquents}

\begin{longtable}{@{}p{0.28\textwidth}p{0.32\textwidth}p{0.32\textwidth}@{}}
\toprule
\textbf{Symptôme} & \textbf{Cause probable} & \textbf{Action} \\
\midrule
\endhead
Redirect login sur page publique & Middleware whitelist & Ajouter route \\
401 Core & Client/API key & Platform Client / secrets \\
kernelIdentityReady=false & Core down / credentials / JWT & Retry identité candidat \\
400 identité incomplète à l'approve & Phase 1 oubliée & JWT candidat kernel-identity \\
403 create org & Pas OWNER / permissions & Parcours onboarding docs Kernel \\
MFA admin bloquante & MFA non activée durablement & Activer MFA compte admin \\
Auth « sur Trust :8090 » & Mauvaise URL & Corriger vers Kernel/Core \\
FleetMan connect fail & Clé AES / plaintext & Vérifier encryption key \\
\bottomrule
\end{longtable}

\section{Dette technique constatée (audit code)}

\begin{enumerate}[leftmargin=*]
  \item \textbf{Dual BFF} --- Next runtime vs Java rollback ; parité non
        garantie route par route.
  \item \textbf{Tests} --- 4 fichiers Vitest ; pas d'E2E.
  \item \textbf{Trust sous-câblé} côté front.
  \item \textbf{Auth asymétrique} --- \api{/branch} et \api{/livreur} publics
        au middleware ; AuthGuard HQ only ; \texttt{PUBLIC\_PREFIXES} mort dans
        AuthGuard.
  \item \textbf{Sessions enrichies coûteuses} --- parcours multi-appels Core
        (N+1) pour branch/livreur session.
  \item \textbf{Envelope hétérogène} --- SUCCESS vs JSON plat.
  \item \textbf{Realtime dual} --- SSE Core vs WS legacy.
  \item \textbf{README frontend partiellement obsolète} --- décrit \api{/} comme
        dashboard (réel : landing ; dashboard = \api{/dashboard}).
  \item \textbf{Messagerie immature}.
  \item \textbf{Nomenclature DTO} encore « tnt-agency » alors que l'upstream
        est Core.
  \item \textbf{Docs ARCHITECTURE} encore centrées DB locale --- à réconcilier
        avec BFF stateless.
\end{enumerate}

\section{Recommandations ingénierie (priorisées)}

\begin{enumerate}
  \item Geler un contrat OpenAPI « Agency BFF » unique (Next) + matrice de
        parité Java ou plan de dépréciation daté.
  \item Homogénéiser envelope + supprimer JSON plat session.
  \item Étendre tests : auth MFA, onboarding approve, offline conflict.
  \item Documenter et optimiser \texttt{auth/*/session} (cache court).
  \item Mettre à jour README frontend (routes réelles).
  \item Décider Trust : client lecture preuves ou purement Core opaque.
\end{enumerate}

\section{Observabilité minimale attendue}

\begin{itemize}
  \item Logs BFF : method, path normalisé, status Core, code erreur envelope
        --- jamais password/api-key/JWT.
  \item Corrélation : \texttt{tenantId}, \texttt{agencyId}, user sub si présent.
  \item Smoke post-deploy : health JWKS + login + tracking public.
\end{itemize}

\chapter{Annexes}

\section{A --- Index des fichiers critiques}

\begin{longtable}{@{}p{0.48\textwidth}p{0.44\textwidth}@{}}
\toprule
\textbf{Chemin} & \textbf{Rôle} \\
\midrule
\endhead
\pathref{frontend/middleware.ts} & Auth edge + SEO headers \\
\pathref{frontend/app/api/agency/[...path]/route.ts} & Entrée BFF \\
\pathref{frontend/lib/server/agency-bff.ts} & Router BFF \\
\pathref{frontend/lib/server/verify-auth.ts} & JWKS / cookies \\
\pathref{frontend/lib/server/fleetman-bff.ts} & Pont FleetMan \\
\pathref{frontend/lib/server/presence-stream.ts} & Relais realtime \\
\pathref{frontend/components/LayoutController.tsx} & Shell vs standalone \\
\pathref{frontend/components/AuthGuard.tsx} & Garde HQ \\
\pathref{frontend/lib/livreur/offlineQueue.ts} & Offline IDB \\
\pathref{frontend/lib/guide/nav.ts} & 17 sections guide \\
\pathref{frontend/lib/seo.ts} & Metadata / sitemap helpers \\
\pathref{frontend/lib/config.ts} & Config publique + GLOBAL\_LINKS \\
\pathref{backend/docs/AUTH\_BFF.md} & Contrat auth \\
\pathref{backend/docs/ONBOARDING\_KERNEL\_FLOW.md} & Onboarding 2 phases \\
\pathref{backend/docs/ADR\_CORE\_FIRST\_MIGRATION.md} & Core-first \\
\pathref{backend/tnt-agency/INTEGRATIONS.md} & Règles BFF Java \\
\bottomrule
\end{longtable}

\section{B --- Séquence login HQ (détaillée)}

\begin{lstlisting}
UI /login
  POST /api/agency/auth/login { email, password }
    BFF -> Core /api/v1/auth/login (+ X-Client-Id/Key)
      alt MFA
        <- 202 mfaToken
        UI MfaChallengeForm
        POST .../mfa/confirm
      else
        Set-Cookie tnt-auth (+ tnt-shared-session?)
  GET /api/agency/auth/session
    resolve agencyId/status
  localStorage persist
  navigate /dashboard
  AuthGuard OK si ACTIVE
\end{lstlisting}

\section{C --- Séquence dépôt QR $\rightarrow$ POD}

\begin{lstlisting}
HQ imprime QR -> /track/deposit?agencyId&branchId
Client POST intake-requests (public)
Agent /accueil approve -> mission PENDING
Antenne/HQ assign -> ASSIGNED
Livreur pickup -> IN_TRANSIT
  (offline: enqueue puis sync)
Livreur deliver + POD -> DELIVERED
Billing/commissions sur mission cloturee
\end{lstlisting}

\section{D --- Guide in-app (produit, pas technique)}

17 slugs (\pathref{lib/guide/nav.ts}) : demarrer, portails, depot-client,
suivi-colis, dashboard, missions, antennes, flotte, hubs, personnel,
vue-antenne, vue-livreur, commissions, facturation, litiges, parametres,
workflow-complet --- avec demos sous \pathref{components/guide/demos/}.

\section{E --- Références externes à maintenir synchrones}

\begin{itemize}
  \item \pathref{ARCHITECTURE\_CORE\_KERNEL\_AGENCY.md} (à réaligner BFF sans DB)
  \item \pathref{CORE\_CONSUMER\_CONTRACT.md}
  \item \pathref{CORE\_CONSUMPTION\_AUDIT.md}
  \item \pathref{KERNEL\_CLIENT\_APPLICATION\_SETUP.md}
  \item Guide Platform Clients Core
\end{itemize}

\section{F --- Comment contribuer à ce document}

\begin{enumerate}
  \item Modifier le chapitre concerné (pas un README parallèle).
  \item Incrémenter le journal des versions dans \pathref{main.tex}.
  \item Si contrat BFF change : mettre à jour ch.~4, 6, 8 et tests Vitest.
  \item Recompiler : \texttt{pdflatex main.tex} (×2).
\end{enumerate}


\end{document}
